\documentclass[aspectratio=169]{beamer}
%--- Configuración del Tema y Colores ---
\usetheme{Madrid}
\usecolortheme{whale}
\setbeamertemplate{navigation symbols}{}
\setbeamertemplate{caption}[numbered]

%--- Paquetes ---
\usepackage[utf8]{inputenc}
\usepackage[spanish]{babel}
\usepackage{graphicx}
\usepackage{booktabs}
\usepackage{tikz}
\usepackage{fontawesome5}
\usetikzlibrary{shapes.geometric, arrows, positioning}

%--- Metadatos del Documento ---
\title[PM2]{Programación Matemática 2}
\subtitle{Clase 5: Iteración Rápida y Revisión de Código con IA}
\author[Luis Alfredo Alvarado]{Prof. Luis Alfredo Alvarado Rodríguez}
\institute[ECFM - USAC]{
    Escuela de Ciencias Físicas y Matemáticas \\
    Universidad de San Carlos de Guatemala
}
\date{Primer Semestre, 2026}

\begin{document}

%--- Diapositiva de Título ---
\begin{frame}
    \titlepage
\end{frame}

%--- Diapositiva: Tabla de Contenidos ---
\begin{frame}{Agenda de Hoy}
    \tableofcontents
\end{frame}

%===========================================================
% SECCIÓN 1: FILOSOFÍA DE ITERACIÓN RÁPIDA
%===========================================================
\section{Filosofía de la Iteración Rápida}

\begin{frame}{El Principio de Iteración Rápida}
    \begin{block}{Concepto Central}
        La \textbf{iteración rápida} consiste en producir versiones funcionales del código lo más pronto posible, obtener feedback inmediato, y refinar incrementalmente.
    \end{block}
    \vspace{0.3cm}
    
    \begin{alertblock}{Mantra del Vibe Coding}
        \textbf{``Hacer que funcione $\rightarrow$ Hacer que funcione bien $\rightarrow$ Hacer que funcione rápido''}
    \end{alertblock}
    \vspace{0.3cm}
    
    \centering
    \textit{No se busca la perfección en el primer intento, sino la mejora continua.}
\end{frame}

\begin{frame}{El Loop de Feedback}
    \centering
    \begin{tikzpicture}[node distance=1.8cm, auto]
        \tikzstyle{step} = [rectangle, rounded corners, minimum width=2.2cm, minimum height=1cm, text centered, draw=black, fill=blue!20, font=\small]
        \tikzstyle{arrow} = [thick,->,>=stealth]
        
        \node (prompt) [step, fill=green!30] {Prompt};
        \node (generate) [step, right=1.2cm of prompt] {Generar};
        \node (execute) [step, right=1.2cm of generate] {Ejecutar};
        \node (evaluate) [step, right=1.2cm of execute] {Evaluar};
        \node (refine) [step, below=1cm of evaluate] {Refinar};
        
        \draw [arrow] (prompt) -- (generate);
        \draw [arrow] (generate) -- (execute);
        \draw [arrow] (execute) -- (evaluate);
        \draw [arrow] (evaluate) -- (refine);
        \draw [arrow] (refine.west) -- ++(-6.5,0) |- (prompt.south);
    \end{tikzpicture}
    
    \vspace{0.5cm}
    \textbf{Objetivo:} Minimizar el tiempo de cada ciclo. Idealmente, segundos a minutos.
\end{frame}

\begin{frame}{Estrategias de Iteración}
    \begin{enumerate}
        \setlength\itemsep{0.8em}
        \item \textbf{Divide y vencerás:} Descomponer problemas grandes en funciones pequeñas e independientes.
        
        \item \textbf{Happy path primero:} Implementar el flujo principal antes de manejar excepciones.
        
        \item \textbf{Hardcoding temporal:} Usar valores fijos para probar lógica, parametrizar después.
        
        \item \textbf{Print debugging:} Agregar prints estratégicos para entender el flujo.
        
        \item \textbf{Fail fast:} Ejecutar frecuentemente para detectar errores temprano.
    \end{enumerate}
\end{frame}

%===========================================================
% SECCIÓN 2: DEBUGGING ASISTIDO POR IA
%===========================================================
\section{Debugging Asistido por IA}

\begin{frame}{El Nuevo Paradigma del Debugging}
    \begin{columns}[t]
        \column{0.48\textwidth}
        \begin{alertblock}{Debugging Tradicional}
            \begin{enumerate}
                \item Leer el mensaje de error
                \item Buscar en Stack Overflow
                \item Probar múltiples soluciones
                \item Agregar prints para rastrear
                \item Usar debugger paso a paso
            \end{enumerate}
        \end{alertblock}
        
        \column{0.48\textwidth}
        \begin{exampleblock}{Debugging con IA}
            \begin{enumerate}
                \item Copiar el error y código
                \item Pegar en el chat de la IA
                \item Recibir explicación y solución
                \item Aplicar el fix sugerido
                \item Iterar si es necesario
            \end{enumerate}
        \end{exampleblock}
    \end{columns}
\end{frame}

\begin{frame}{Técnicas de Debugging con IA}
    \begin{enumerate}
        \setlength\itemsep{0.7em}
        \item \textbf{Explicación de errores:} ``Explica este error y por qué ocurre''
        
        \item \textbf{Rastreo de flujo:} ``Muéstrame paso a paso qué hace este código con input X''
        
        \item \textbf{Identificación de edge cases:} ``¿Qué inputs podrían causar problemas?''
        
        \item \textbf{Comparación:} ``¿Por qué este código funciona pero este otro no?''
        
        \item \textbf{Generación de tests:} ``Genera casos de prueba que podrían fallar''
        
        \item \textbf{Rubber duck con esteroides:} Explicar el problema a la IA para clarificar el pensamiento.
    \end{enumerate}
\end{frame}

%===========================================================
% SECCIÓN 3: REFACTORIZACIÓN CON IA
%===========================================================
\section{Refactorización con IA}

\begin{frame}{¿Qué es la Refactorización?}
    \begin{block}{Definición}
        La \textbf{refactorización} es el proceso de reestructurar código existente sin cambiar su comportamiento externo, mejorando su estructura interna.
    \end{block}
    \vspace{0.3cm}
    
    \textbf{Objetivos:}
    \begin{itemize}
        \item Mejorar la legibilidad
        \item Reducir la complejidad
        \item Eliminar duplicación
        \item Mejorar el rendimiento
        \item Facilitar el mantenimiento
    \end{itemize}
    \vspace{0.3cm}
    
    \centering
    \textit{``Los buenos programadores escriben código que los humanos entienden.''} — Martin Fowler
\end{frame}

\begin{frame}{Tipos de Refactorización con IA}
    \begin{columns}[t]
        \column{0.48\textwidth}
        \begin{block}{Estructural}
            \begin{itemize}
                \item Extraer funciones
                \item Simplificar condicionales
                \item Eliminar código muerto
                \item Renombrar variables
            \end{itemize}
        \end{block}
        
        \column{0.48\textwidth}
        \begin{block}{Rendimiento}
            \begin{itemize}
                \item Optimizar algoritmos
                \item Usar estructuras de datos correctas
                \item Eliminar operaciones redundantes
                \item Aplicar patrones eficientes
            \end{itemize}
        \end{block}
    \end{columns}
    
    \vspace{0.5cm}
    \begin{exampleblock}{Prompts Útiles}
        ``Refactorizar siguiendo SRP'' | ``Optimizar de O(n²) a O(n)'' | ``Convertir a list comprehension''
    \end{exampleblock}
\end{frame}

%===========================================================
% SECCIÓN 4: CODE REVIEW CON IA
%===========================================================
\section{Code Review con IA}

\begin{frame}{Code Review Automatizado}
    \begin{block}{¿Qué puede revisar la IA?}
        \begin{itemize}
            \item \textbf{Bugs potenciales:} Null references, off-by-one, race conditions
            \item \textbf{Seguridad:} SQL injection, XSS, manejo de secretos
            \item \textbf{Estilo:} PEP 8, convenciones del proyecto
            \item \textbf{Rendimiento:} Algoritmos ineficientes, memory leaks
            \item \textbf{Mantenibilidad:} Código duplicado, funciones muy largas
        \end{itemize}
    \end{block}
    
    \vspace{0.3cm}
    \centering
    \textit{La IA no reemplaza el code review humano, lo complementa.}
\end{frame}

\begin{frame}{Checklist de Code Review}
    \begin{columns}[t]
        \column{0.48\textwidth}
        \begin{block}{Funcionalidad}
            \begin{itemize}
                \item[\faCheckSquare] ¿Hace lo que debe?
                \item[\faCheckSquare] ¿Maneja edge cases?
                \item[\faCheckSquare] ¿Errores manejados?
                \item[\faCheckSquare] ¿Inputs validados?
            \end{itemize}
        \end{block}
        
        \begin{block}{Seguridad}
            \begin{itemize}
                \item[\faCheckSquare] ¿Inputs sanitizados?
                \item[\faCheckSquare] ¿Secrets seguros?
                \item[\faCheckSquare] ¿Sin vulnerabilidades?
            \end{itemize}
        \end{block}
        
        \column{0.48\textwidth}
        \begin{block}{Rendimiento}
            \begin{itemize}
                \item[\faCheckSquare] ¿Complejidad adecuada?
                \item[\faCheckSquare] ¿Operaciones necesarias?
                \item[\faCheckSquare] ¿Estructuras correctas?
            \end{itemize}
        \end{block}
        
        \begin{block}{Mantenibilidad}
            \begin{itemize}
                \item[\faCheckSquare] ¿Código legible?
                \item[\faCheckSquare] ¿Responsabilidad única?
                \item[\faCheckSquare] ¿Tests incluidos?
            \end{itemize}
        \end{block}
    \end{columns}
\end{frame}



%===========================================================
% SECCIÓN 6: EJERCICIOS
%===========================================================
\section{Ejercicios Propuestos}

\begin{frame}{Ejercicios de Iteración y Debugging}
    \begin{enumerate}
        \setlength\itemsep{0.8em}
        \item \textbf{Debugging:} Corregir una función con 3 bugs usando IA para identificarlos.
        
        \item \textbf{Refactorización:} Tomar una función de 50+ líneas y dividirla siguiendo SRP.
        
        \item \textbf{Optimización:} Convertir una función $O(n^2)$ a $O(n)$ con ayuda de IA.
        
        \item \textbf{Code Review:} Realizar review completo de un módulo, identificando 5 mejoras.

    \end{enumerate}
\end{frame}

%--- Diapositiva Final ---
\begin{frame}
    \centering
    \Huge \textbf{¿Preguntas?}
    
    \vspace{0.8cm}
    \large
    \textbf{Próxima clase:} Ingeniería de Software: Git
    
    \vspace{0.5cm}
    \normalsize
    Prof. Luis Alfredo Alvarado Rodríguez\\
    \small Escuela de Ciencias Físicas y Matemáticas
    
    \vspace{0.5cm}
    \footnotesize
    \textit{``Genera rápido, itera hasta que funcione.''}
\end{frame}

\end{document}